% GR Chain Bridge — covariant-tetrad revision, 2026-07-15.
\documentclass[11pt,a4paper]{article}
\usepackage{amsmath,amssymb}
\usepackage[most]{tcolorbox}
\usepackage[margin=2.5cm]{geometry}
\begin{document}
\section{UBT to GR: current canonical bridge}

The canonical local chain is
\[
 \Theta(q,\tau)
 \longrightarrow E_\mu=\mathcal N_0^{-1/2}D_\mu\Theta
 \longrightarrow
 \frac12(E_\mu^\sharp E_\nu+E_\nu^\sharp E_\mu)
 =g_{\mu\nu}\mathbf1.
\]
The metric identity and local rank-ten theorem are proved in
\texttt{canonical/gr\_closure/covariant\_tetrad\_rank\_theorem.tex}.
The same four $E_\mu$ admit the injective block lift
\[
 \Gamma_\mu=\mathcal C(E_\mu),\qquad
 \tfrac12\{\Gamma_\mu,\Gamma_\nu\}=g_{\mu\nu}I_4,
\]
so the generalized-Dirac operator is an operator completion of the canonical
relation, not a second metric construction.  The exact lift is proved in
\texttt{canonical/geometry/biquaternion\_dirac\_lift.tex}.

\section{Connection and curvature}
The frame connection $\Omega_\mu$ and affine connection
$\Gamma^\rho{}_{\mu\nu}$ are related by the tetrad postulate
\[
 \partial_\mu E_\nu-\Gamma^\rho{}_{\mu\nu}E_\rho
 +\rho_{\mathrm{vec}*}(\Omega_\mu)E_\nu=0.
\]
They are not related by a componentwise real projection.  Curvature is
\[
 [D_\mu,D_\nu]
 =\partial_\mu\Omega_\nu-\partial_\nu\Omega_\mu
 +[\Omega_\mu,\Omega_\nu].
\]

\section{Proof status}
\begin{tcolorbox}[colback=green!5!white,colframe=green!60!black,title=Closed]
The local tetrad-to-metric map has rank ten and reproduces every local Lorentz
metric up to six local Lorentz frame directions.
\end{tcolorbox}

\begin{tcolorbox}[colback=yellow!7!white,colframe=orange!70!black,title=Open]
The UBT action must still derive or uniquely eliminate $\Omega_\mu$, preserve
the Lorentz slice, solve the integrability problem $D_\mu\Theta=E_\mu$, and
derive Einstein dynamics.  The complete GR chain is therefore
\textbf{DERIVED WITH ASSUMPTIONS / OPEN DYNAMICAL BRIDGE}, not fully closed.
\end{tcolorbox}

The former compact-$\psi$ fiber-average route is exploratory and noncanonical.
\end{document}
